\subsubsubsection{NTT 原根表}

常用的 NTT 模数及其原根如下（$r\cdot 2^k+1$ 形式）：

\begin{table}[htbp]
\centering
\arrayrulecolor{black}
\setlength{\arrayrulewidth}{0.8pt}
\begin{tabular}{|c|c|c|c|}
\hline
$r\cdot 2^k+1$ & $r$ & $k$ & 原根 $g$ \\
\hline
\textbf{3} & 1 & 1 & 2 \\
\textbf{5} & 1 & 2 & 2 \\
\textbf{17} & 1 & 4 & 3 \\
\textbf{97} & 3 & 5 & 5 \\
\textbf{193} & 3 & 6 & 5 \\
\textbf{257} & 1 & 8 & 3 \\
\textbf{7681} & 15 & 9 & 17 \\
\textbf{12289} & 3 & 12 & 11 \\
\textbf{40961} & 5 & 13 & 3 \\
\textbf{65537} & 1 & 16 & 3 \\
\textbf{786433} & 3 & 18 & 10 \\
\textbf{5767169} & 11 & 19 & 3 \\
\textbf{7340033} & 7 & 20 & 3 \\
\textbf{23068673} & 11 & 21 & 3 \\
\textbf{104857601} & 25 & 22 & 3 \\
\textbf{167772161} & 5 & 25 & 3 \\
\textbf{469762049} & 7 & 26 & 3 \\
\cellcolor[gray]{0.8}{\textbf{998244353}} & 119 & 23 & 3 \\
\cellcolor[gray]{0.8}{\textbf{1004535809}} & 479 & 21 & 3 \\
\textbf{2013265921} & 15 & 27 & 31 \\
\cellcolor[gray]{0.8}{2281701377} & 17 & 27 & 3 \\
3221225473 & 3 & 30 & 5 \\
75161927681 & 35 & 31 & 3 \\
77309411329 & 9 & 33 & 7 \\
206158430209 & 3 & 36 & 22 \\
2061584302081 & 15 & 37 & 7 \\
2748779069441 & 5 & 39 & 3 \\
6597069766657 & 3 & 41 & 5 \\
39582418599937 & 9 & 42 & 5 \\
79164837199873 & 9 & 43 & 5 \\
263882790666241 & 15 & 44 & 7 \\
1231453023109121 & 35 & 45 & 3 \\
1337006139375617 & 19 & 46 & 3 \\
3799912185593857 & 27 & 47 & 5 \\
4222124650659841 & 15 & 48 & 19 \\
7881299347898369 & 7 & 50 & 6 \\
31525197391593473 & 7 & 52 & 3 \\
180143985094819841 & 5 & 55 & 6 \\
1945555039024054273 & 27 & 56 & 5 \\
4179340454199820289 & 29 & 57 & 3 \\
\hline
\end{tabular}
\caption{NTT 常用模数与原根（加粗为 int 范围，灰色背景为常用模数）}
\end{table}

\par
\textbf{说明：}

\begin{itemize}
    \item 加粗的数字（如 \textbf{3}、\textbf{998244353} 等）为 \texttt{int} 类型可存储的模数。
    \item 带有灰色背景的 $998244353$、$1004535809$、$2281701377$ 是竞赛中最常用的三个模数。
    \item $2281701377 = 17\times 2^{27}+1$，平方不会爆 \texttt{long long}。
    \item $1004535809 = 479\times 2^{21}+1$，加法不会爆 \texttt{int}。
\end{itemize}